(a) Regular functions glue uniquely, and vanishing on $Y$ is a local condition. Therefore sections in the prescribed ideals glue uniquely to a section in the same ideal, so $\mathscr I_Y$ is a sheaf of ideals in $\mathcal O_X$.

(b) Restricting functions gives $\mathcal O_X\to i_*\mathcal O_Y$ with kernel $\mathscr I_Y$. At $P\in Y$, choose an affine neighborhood with ring $A$ and let $I$ define $Y$ there. The map on stalks is $A_{\mathfrak m_P}\to(A/I)_{\mathfrak m_P/I}$, which is surjective with kernel $IA_{\mathfrak m_P}$. Outside $Y$ the target stalk is zero. Thus the morphism is surjective, and \exref{II.1.6} gives $\mathcal O_X/\mathscr I_Y\cong i_*\mathcal O_Y$.

(c) Evaluation at $P,Q$ gives $\mathcal O_X\to i_*\mathcal O_P\oplus i_*\mathcal O_Q$. It is surjective on stalks, and its kernel consists of functions vanishing at both points. On global sections it is the diagonal map $k\to k\oplus k$, since $\mathcal O(\mathbb{P}^1)=k$, so it is not surjective.

(d) Every regular function on a nonempty open subset is a rational function, giving $\mathcal O\to\mathscr K$. Send a rational function to its classes in $K/\mathcal O_P$ at all points. It has only finitely many poles, so these classes define a section of $\bigoplus_P i_P(K/\mathcal O_P)$. This is a morphism of sheaves, and its stalk map at $Q$ is the quotient $K\to K/\mathcal O_Q$. Therefore its kernel is $\mathcal O$ and it is surjective, yielding $\mathscr K/\mathcal O\cong\bigoplus_P i_P(K/\mathcal O_P)$.

(e) Since $\mathbb{P}^1$ is quasi-compact, a global section of this direct sum has only finitely many nonzero components. In the coordinate $t$, every class at a finite point $a$ has a representative $\sum_{j\geq1}c_{a,j}(t-a)^{-j}$ with finitely many terms; a class at infinity has a representative $\sum_{j\geq1}c_{\infty,j}t^j$. The sum of these representatives has exactly the prescribed class at every point, since each summand is regular away from its indicated point. Thus $K\to\Gamma(X,\mathscr K/\mathcal O)$ is surjective. Its kernel consists of rational functions regular everywhere, namely $k=\Gamma(X,\mathcal O)$, proving exactness.
